\subsection*{CU-24: Abandonar la partida}
\addcontentsline{toc}{subsection}{CU-24: Abandonar la partida}
\label{cu-24}

\begin{fichacu}{CU-24}{Abandonar la partida}
  \crudFila{Responsable}{Ángel Gabriel Aguilar Hernández}
  \crudFila{Actor principal}{Jugador}
  \crudFila{Actores de sistema}{Servidor de partidas, Base de datos}
  \crudFila{Prototipos}{8e, 6f, 7g}
  \crudFila{Objetivo}{Permitir al Jugador salir definitivamente de una partida en curso, exigiéndole escribir la confirmación, y resolver también el abandono que se produce cuando un jugador desconectado no vuelve y su rival reclama la victoria.}
  \crudFila{Disparador}{El Jugador intenta salir de la \texttt{GUI\_\allowbreak{}Match} —cerrar la pantalla, volver al menú o cerrar el juego— con la partida en curso; o el rival de un jugador desconectado selecciona ``reclamar victoria''.}
  \textbf{Precondiciones} & \cuCelda{\begin{cureglas}
      \item \textbf{\mbox{PRE-01}} El Jugador tiene una \textit{Participacion} sin terminar en una \textit{Partida} con \texttt{estado} igual a \texttt{EN\_\allowbreak{}CURSO}.
      \item \textbf{\mbox{PRE-02}} El Servidor de partidas está en ejecución y tiene conexión disponible con la Base de datos.
    \end{cureglas}} \\ \hline
  \textbf{Flujo normal} & \cuCelda{\begin{cuenum}[start=1]
      \item El Jugador intenta salir de la \texttt{GUI\_\allowbreak{}Match}.
      \item El SJM despliega la \texttt{GUI\_\allowbreak{}MessageConfirm} de nivel más alto, que explica que abandonar cuenta como derrota con su pérdida de elo, que el peón se retira pero los muros se quedan, y pide escribir la palabra de confirmación indicada, con las opciones ``Abandonar'' y ``Volver a la partida'' (\mbox{RN-02}). \textit{(Ver \mbox{FA-01}, \mbox{FA-02}, \mbox{FA-08})}
      \item El Jugador escribe la palabra y selecciona ``Abandonar''.
      \item El SJM verifica que la palabra coincida y envía el mensaje \texttt{LEAVE\_\allowbreak{}MATCH\_\allowbreak{}REQUEST} con el token de sesión. \textit{(Ver \mbox{EX-02}, \mbox{EX-03})}
      \item El Servidor de partidas verifica que el token corresponda a una \textit{Sesion} abierta y a una \textit{Participacion} sin terminar de esa partida. \textit{(Ver \mbox{FA-03})}
      \item El Servidor de partidas verifica que la partida siga \texttt{EN\_\allowbreak{}CURSO}. \textit{(Ver \mbox{FA-04})}
      \item El Servidor de partidas marca como \texttt{CADUCADA} cualquier \textit{OfertaTablas} pendiente.
      \item El Servidor de partidas incrementa \texttt{abandonos} en la \textit{EstadisticaModo} del Jugador para el modo de la partida (\mbox{RN-05}). \textit{(Ver \mbox{EX-01})}
    \end{cuenum}} \\
   & \cuCelda{\begin{cuenum}[start=9]
      \item Si la partida es de dos jugadores, el flujo continúa en el \mbox{CU-25} Finalizar la partida con la forma de término \texttt{ABANDONO} y el rival como ganador. \textit{(Ver \mbox{FA-05})}
      \item El SJM del Jugador despliega el menú principal con el resultado de la partida.
      \item El SJM del rival muestra que el Jugador abandonó y despliega la pantalla de fin de partida.
      \item Fin del caso de uso.
    \end{cuenum}} \\ \hline
  \textbf{Flujos alternos} & \cuCelda{\textbf{\mbox{FA-01}: El Jugador vuelve a la partida}\par
    \begin{cuenum}[start=1]
      \item El Jugador selecciona ``Volver a la partida''.
      \item El SJM cierra la confirmación y devuelve el foco al tablero.
      \item Fin del caso de uso. El reloj siguió corriendo durante la confirmación.
    \end{cuenum}} \\
   & \cuCelda{\textbf{\mbox{FA-02}: El Jugador cierra el juego sin confirmar}\par
    \begin{cuenum}[start=1]
      \item El SJM no puede mostrar la confirmación porque el sistema operativo cierra la aplicación.
      \item El Servidor de partidas detecta la pérdida de conexión y \textbf{no} lo trata como abandono: lo trata como desconexión, conforme al \mbox{CU-26} (\mbox{RN-03}).
      \item Fin del caso de uso.
    \end{cuenum}} \\
   & \cuCelda{\textbf{\mbox{FA-03}: La sesión o la participación no son válidas}\par
    \begin{cuenum}[start=1]
      \item El Servidor de partidas responde con \texttt{SESSION\_\allowbreak{}INVALID} o \texttt{NOT\_\allowbreak{}A\_\allowbreak{}PARTICIPANT}.
      \item El SJM despliega la \texttt{GUI\_\allowbreak{}Login} o el menú principal, según corresponda.
      \item Fin del caso de uso.
    \end{cuenum}} \\
   & \cuCelda{\textbf{\mbox{FA-04}: La partida ya terminó}\par
    \begin{cuenum}[start=1]
      \item El Servidor de partidas responde con \texttt{MATCH\_\allowbreak{}ALREADY\_\allowbreak{}FINISHED} y el resultado real.
      \item El SJM sale de la partida sin aplicar el abandono y muestra ese resultado.
      \item Fin del caso de uso.
    \end{cuenum}} \\
   & \cuCelda{\textbf{\mbox{FA-05}: La partida es de cuatro jugadores}\par
    \begin{cuenum}[start=1]
      \item El Servidor de partidas cierra solo la \textit{Participacion} del Jugador con resultado perdido, forma de término \texttt{ABANDONO} y la última posición entre los activos, y retira su peón del tablero.
      \item Los muros que había colocado permanecen en el tablero y siguen bloqueando (\mbox{RN-04}).
      \item Si quedan al menos dos jugadores activos, la partida sigue y el turno salta al Jugador; si queda uno solo, el flujo continúa en el \mbox{CU-25}.
      \item Los SJM de los demás muestran el aviso de abandono y la fila del jugador tachada.
      \item Fin del caso de uso.
    \end{cuenum}} \\
   & \cuCelda{\textbf{\mbox{FA-06}: El rival reclama la victoria sobre un jugador desconectado}\par
    \begin{cuenum}[start=1]
      \item El SJM del rival, que lleva sesenta segundos viendo la espera de la \texttt{GUI\_\allowbreak{}OpponentDisconnected}, habilita ``Reclamar victoria'' (\mbox{RN-06}).
      \item El rival la selecciona y su SJM envía \texttt{CLAIM\_\allowbreak{}VICTORY\_\allowbreak{}REQUEST}.
      \item El Servidor de partidas verifica que la \textit{Participacion} del jugador desconectado tenga \texttt{conectado} en falso y \texttt{fecha\_\allowbreak{}desconexion} anterior a sesenta segundos. \textit{(Ver \mbox{FA-07})}
      \item El Servidor de partidas aplica al desconectado los pasos 7 y 8 del FN.
      \item El flujo continúa en el paso 9 del FN, con el jugador desconectado como el que abandona.
    \end{cuenum}} \\
   & \cuCelda{\textbf{\mbox{FA-07}: El jugador desconectado volvió antes de reclamar}\par
    \begin{cuenum}[start=1]
      \item El Servidor de partidas encuentra la \textit{Participacion} con \texttt{conectado} en verdadero.
      \item El Servidor de partidas responde con \texttt{OPPONENT\_\allowbreak{}RECONNECTED} y no aplica nada.
      \item El SJM del rival retira la espera y la partida sigue.
      \item Fin del caso de uso.
    \end{cuenum}} \\
   & \cuCelda{\textbf{\mbox{FA-08}: El Jugador abandona una partida contra la IA}\par
    \begin{cuenum}[start=1]
      \item El SJM muestra una confirmación sencilla, sin escribir la palabra, porque no hay elo ni rival en juego.
      \item El Servidor de partidas cierra la partida con la forma de término \texttt{ABANDONO} conforme al \mbox{FA-01} del \mbox{CU-25}, sin incrementar \texttt{abandonos} (\mbox{RN-05}).
      \item Fin del caso de uso.
    \end{cuenum}} \\ \hline
  \textbf{Excepciones} & \cuCelda{\textbf{\mbox{EX-01}: Fallo al escribir en la base de datos}\par
    \begin{cuenum}[start=1]
      \item El Servidor de partidas revierte lo escrito y registra la excepción con un identificador de incidencia.
      \item El Servidor de partidas responde con \texttt{SERVER\_\allowbreak{}ERROR} sin aplicar el abandono.
      \item El SJM informa del fallo y deja al Jugador en la partida, que sigue en curso.
      \item Fin del caso de uso.
    \end{cuenum}} \\
   & \cuCelda{\textbf{\mbox{EX-02}: Se agota el tiempo de espera de la respuesta}\par
    \begin{cuenum}[start=1]
      \item El SJM sale igualmente de la pantalla de partida, porque el Jugador ya confirmó que quiere irse.
      \item Si el abandono no llegó al servidor, este tratará la salida como desconexión conforme al \mbox{FA-02}, y el rival podrá reclamar la victoria.
      \item Fin del caso de uso.
    \end{cuenum}} \\
   & \cuCelda{\textbf{\mbox{EX-03}: La conexión se interrumpe durante el intercambio}\par
    \begin{cuenum}[start=1]
      \item El flujo continúa en el paso 2 del \mbox{EX-02}.
    \end{cuenum}} \\ \hline
  \textbf{Postcondiciones} & \cuCelda{\begin{cureglas}
      \item \textbf{\mbox{POST-01}} Si el caso termina por el FN o el \mbox{FA-06}, la partida está cerrada con la forma de término \texttt{ABANDONO}, el que abandona perdió y su rival ganó.
      \item \textbf{\mbox{POST-02}} Si el caso termina por el FN o el \mbox{FA-06} en una partida en línea, la \textit{EstadisticaModo} del que abandona tiene \texttt{abandonos} incrementado.
      \item \textbf{\mbox{POST-03}} Si el caso termina por el \mbox{FA-05}, solo la \textit{Participacion} del que abandona está cerrada y sus muros siguen en el tablero.
      \item \textbf{\mbox{POST-04}} Si el caso termina por el \mbox{FA-01}, el \mbox{FA-02}, el \mbox{FA-04} o el \mbox{FA-07}, no se aplicó ningún abandono.
    \end{cureglas}} \\ \hline
  \textbf{Reglas de negocio} & \cuCelda{\begin{cureglas}
      \item \textbf{\mbox{RN-01}} Abandonar cuenta como una derrota a todos los efectos: elo, contadores y racha.
      \item \textbf{\mbox{RN-02}} Abandonar exige escribir la palabra de confirmación. Es la confirmación más fuerte de las tres que frenan una partida, porque deja al rival sin partida.
      \item \textbf{\mbox{RN-03}} Cerrar el juego o perder la conexión no es abandonar. Es una desconexión, y la partida espera al jugador mientras su reloj corre (\mbox{CU-26}).
      \item \textbf{\mbox{RN-04}} El peón del que abandona se retira del tablero; sus muros permanecen (\mbox{RN-10} del \mbox{CU-25}).
      \item \textbf{\mbox{RN-05}} Los abandonos en partidas en línea se cuentan por modo en \texttt{abandonos} de la \textit{EstadisticaModo}. El contador no se muestra en el perfil, pero está disponible para la moderación.
      \item \textbf{\mbox{RN-06}} El rival de un jugador desconectado puede reclamar la victoria a partir de los sesenta segundos de desconexión. Antes de eso debe esperar, porque la mayoría de las desconexiones son breves.
      \item \textbf{\mbox{RN-07}} Reclamar la victoria convierte la desconexión en abandono, con las mismas consecuencias que un abandono voluntario.
    \end{cureglas}} \\ \hline
  \textbf{Observación} & \cuCelda{\textit{Sobre por qué abandonar y rendirse son casos distintos.}\par
    El resultado de ambos es una derrota, y podría parecer que sobra uno. Se separan porque su intención y su riesgo son distintos: rendirse es un gesto dentro de la partida, que el jugador hace mirando el tablero; abandonar es salir de ella, muchas veces sin mirar. Por eso confirman con distinta fuerza, y por eso solo el abandono se cuenta aparte: un jugador que se rinde mucho juega mal, pero uno que abandona mucho deja a otros sin partida, y eso sí es asunto de moderación.} \\ \hline
  \textbf{Datos que toca} & \cuCelda{\begin{cudatos}
      \item \textbf{\textit{Sesion}, \textit{Participacion}:} lee por token \textit{(paso 5)}
      \item \textbf{\textit{Partida}, \textit{Modo}:} lee \texttt{estado}, \texttt{tipo} y número de jugadores \textit{(pasos 6, 9)}
      \item \textbf{\textit{OfertaTablas}:} marca las pendientes como \texttt{CADUCADA} \textit{(paso 7)}
      \item \textbf{\textit{EstadisticaModo}:} incrementa \texttt{abandonos} \textit{(paso 8)}
      \item \textbf{\textit{Participacion}:} lee \texttt{conectado} y \texttt{fecha\_\allowbreak{}desconexion} del rival; escribe resultado y posición en el \mbox{FA-05} \textit{(pasos \mbox{FA-05}, \mbox{FA-06})}
    \end{cudatos}} \\ \hline
\end{fichacu}
\clearpage
